\begin{abstract}
Sparse autoencoders (SAEs) are widely used to explain what large language models represent, which presumes that an
SAE is a faithful stand-in for the activations it replaces. We audit that presumption at all twelve layers of GPT-2
Small for three SAE architectures (ReLU, TopK and a sigmoid-surrogate JumpReLU; 108 SAEs at matched sparsity, three
seeds each), separating four levels of faithfulness: reconstruction, task-information preservation, behavioural
faithfulness and causal faithfulness. Reconstruction degrades with depth (OWT test explained variance 0.93--0.95 at
layers 0--2, 0.76--0.82 at layer 11). Within a layer, reconstruction error predicts how much splicing the SAE damages
the model's predictions, but not how much truth information is lost. The apparent collapse of a frozen truth probe
on SAE reconstructions is largely reproduced by random noise of the same norm; re-trained readers show a smaller but
real loss, and at equal reconstruction error the SAEs never preserve task information better than PCA. Using an
equal-length persona-instruction task built to remove a length shortcut, we locate where GPT-2 carries the
instruction's effect and decompose it: the SAE-reconstructed activation difference restores 93--97\,\% of the effect
at the layer-1 persona token but only 62--68\,\% at the layer-8 and layer-11 final token, with the remainder in the
reconstruction error or in an interaction between reconstruction and error. Short fine-tuning at layers 3, 6, 8 and 9
separates these levels again. Reconstruction fine-tuning on truth statements cuts their reconstruction error by
61--85\,\% and lowers the truth gap measured by a probe frozen on original activations, meeting its pre-registered
criterion in 4 of 9 layer--architecture cells beyond layer 6; re-trained probes, however, show no consistent gain,
so we do not claim that truth information is recovered. Adding an output-KL term improves behavioural faithfulness
in 5 of 9 creditable cells (in the other 4 the SAE's sparsity drifts beyond the pre-registered guard) and never
improves truth preservation. All results except the reconstruction table are \evlabel{validation} findings, not confirmed results;
a single pre-registered held-out evaluation is pending.
\end{abstract}

\pending{Abstract to be revised after E1--E3 and the final held-out evaluation (M5). Numbers above are validation-stage unless marked \evlabel{test}.}
